\section{Faithful two-face shaping: reachable equilibria of a modal source}\label{sec:faithful}

The finite-array inverse of the preceding sections fits the required load at
the target geometry. A numerical audit of such a fit (study S06 of the
research record) found an equilibrium that matches the target only on its
design grid, moves by micrometres under part-per-thousand changes of wave
parameters, and is not the state reached from the unforced liquid: under the
same command a nearly mirror-image equilibrium attracts the flat start. This
section defines the shaping problem by the equilibrium that the physical
start-up actually reaches, restricts the source to a low-order continuous
modal space, and states the conditions under which the target is formed
faithfully. The statements are conditional on the declared quiescent,
axisymmetric, creeping-flow reduction; non-axisymmetric disturbances, inertia,
streaming and heating remain the separate questions of \cref{sec:stability}
and \cref{sec:validity}.

\subsection{Modal source space}
On each source boundary $\Sigma_b$ (lower window, upper window, cylinder wall)
let the source velocity be a finite combination of fixed real functions,
\begin{equation}
 u(\bm x)=\sum_{k=1}^{m}c_k\,\psi_k(\bm x),\qquad c\in\C^m .
 \label{eq:modal-source}
\end{equation}
On a window of radius $R$ we use the axisymmetric Neumann eigenfunctions
$\psi(r)=J_0(j'_{0,k}r/R)$, with $j'_{0,k}$ the zeros of $J_0'$ including
$j'_{0,1}=0$; on the wall, of height $H$, the cosines $\psi(z)=\cos(k\pi(z-z_0)/H)$.
Each $\psi_k$ is normalized to unit mean square on its boundary, so that
$\lvert c_k\rvert$ is a velocity amplitude in m/s. Low order is a deliberate
restriction: the source has no structure finer than the half-wavelength of its
highest mode, which excludes the cancelling fine-scale fields an
underdetermined fit can use. Any finite partition of $\Sigma_b$ into ideal
regions with region averages of $u$ is a quadrature of
\cref{eq:modal-source}; the region count is then a numerical parameter whose
convergence must be shown, not an apparatus property.

Writing the region values as $g=\Phi c$, every generalized force of
\cref{eq:continuous-quadratic} becomes a Hermitian form in the coefficients,
\begin{equation}
 F_i(q;c)=c^{\mathsf H}\tilde{\bm K}_i(q)\,c,\qquad
 \tilde{\bm K}_i=\Phi^{\mathsf H}\bm K_i\Phi ,
 \label{eq:modal-force}
\end{equation}
with derivative $\partial F_i/\partial(\Rea c,\Ima c)=2(\Rea\tilde{\bm K}_ic,
\Ima\tilde{\bm K}_ic)$. The global phase $c\mapsto e^{\mathrm i\vartheta}c$
leaves every force unchanged; one coefficient phase is fixed.

\subsection{Held-command equilibria and the start-up branch}
Let $q$ collect both interface coefficient vectors in a space that already
enforces the pinned rims and both phase volumes. A held-command equilibrium
satisfies
\begin{equation}
 \mathcal R(q;c)=F_{\mathrm{ac}}(q;c)-\nabla E(q)=0,
 \label{eq:held-equilibrium}
\end{equation}
where $E$ is the capillary--gravity energy, $K=\nabla^2E$ and
$A=\partial F_{\mathrm{ac}}/\partial q$ at fixed $c$. Equation
\cref{eq:held-equilibrium} generally has several solutions. Start-up applies
$s(t)c$ with a ramp $s$ rising from $0$ to $1$. Because every force is
quadratic in the command, the instantaneous load is $s^2F_{\mathrm{ac}}(q;c)$.
Define the start-up branch $q_{\mathcal B}(s)$ by
\begin{equation}
 s^2F_{\mathrm{ac}}(q_{\mathcal B};c)=\nabla E(q_{\mathcal B}),\qquad
 q_{\mathcal B}(0)=q_0,
 \label{eq:startup-branch}
\end{equation}
with $q_0$ the acoustics-off equilibrium $\nabla E(q_0)=0$ at the declared
phase volumes. This start is generally curved: with pinned rims, the fill
volume of each phase fixes the displaced volume of its interface, and gravity
and the Laplace pressure balance then select a meniscus. Along the branch,
\begin{equation}
 \bigl(K-s^2A\bigr)\frac{\dd q_{\mathcal B}}{\dd s}=2s\,F_{\mathrm{ac}} .
 \label{eq:branch-tangent}
\end{equation}

\begin{proposition}[Slow start-up follows the branch]\label{prop:adiabatic}
Assume the creeping-flow reduction $\dot q=M(q)\bigl(s^2F_{\mathrm{ac}}
(q;c)-\nabla E(q)\bigr)$ with a symmetric positive definite mobility $M$, and
suppose that for all $s\in[0,1]$ the branch exists and every eigenvalue of
$-M(K-s^2A)$ has real part at most $-\lambda<0$. Then for a ramp with
$\lvert\dot s\rvert\leq\epsilon$ and $\epsilon$ small relative to $\lambda$ and
the branch curvature, the trajectory from $q_0$ remains within
$O(\epsilon/\lambda)$ of $q_{\mathcal B}(s(t))$ and converges to
$q_{\mathcal B}(1)$ after the ramp ends.
\end{proposition}

This is the standard slow-manifold estimate for a uniformly stable family of
equilibria of a slowly varying vector field \cite{fenichel1979}; the deviation obeys a linear system with decay rate at
least $\lambda$ forced by $\dot s\,\partial_s q_{\mathcal B}$. If $K-s^2A$
becomes singular at some $s_f<1$ (a fold), the branch ends there and the
trajectory jumps to another attractor: the start-up then forms a different
state, as in S06. The reachable equilibrium is
$q_{\mathcal R}(c)=q_{\mathcal B}(1)$ when the hypotheses hold. The branch is a
property of the physical start-up, not continuation in the target; whether a
computed trajectory actually follows it is checked independently by
integrating the forward model from $q_0$.

\subsection{Faithful inverse}
Let $H$ sample both pupils, $q^\star$ be the target, and
$\rho(c)=H(q_{\mathcal R}(c)-q^\star)$. Only the pupils are prescribed: the
non-optical annuli take whatever equilibrium shape the command produces,
subject to pinning and volume. Differentiating \cref{eq:held-equilibrium} at
$s=1$,
\begin{equation}
 \frac{\partial\rho}{\partial(\Rea c,\Ima c)}
 =H\,(K-A)^{-1}\,2\bigl(\Rea\tilde{\bm K}c,\ \Ima\tilde{\bm K}c\bigr),
 \label{eq:faithful-jacobian}
\end{equation}
with $\tilde{\bm K}c$ the matrix of rows $\tilde{\bm K}_ic$. The coupled
operator $(K-A)^{-1}$ replaces the frozen $K^{-1}$ of the earlier inverse;
the two differ by the acoustic shape stiffness, which at S06 amplified
force errors about sixfold. The faithful inverse is
\begin{equation}
 \min_{c}\ \tfrac12\norm{W_h\rho(c)}^2
 +\tfrac12\sum_{\theta}\norm{W_h H\,\delta q_\theta\,\Delta\theta}^2
 \quad\text{subject to}\quad
 \mathcal P(c)\leq P_{\max},\ \ \max_j\lvert p_j(c)\rvert\leq p_{\max},
 \label{eq:faithful-problem}
\end{equation}
where $W_h$ scales heights by the accuracy goal, $\delta q_\theta
=-(K-A)^{-1}\partial_\theta\mathcal R$ is the first-order equilibrium shift per
unit of an uncertain parameter $\theta$ (temperature through the sound speeds,
interfacial tension, frequency) with declared tolerance $\Delta\theta$, the
source power $\mathcal P(c)=c^{\mathsf H}\bm W_{\mathcal P}c$ is a Hermitian
form and $p_j=\bm s_j^{\mathsf T}c$ are sampled pressures. Both caps are
convex in $c$. The robustness term is necessary because a residual that
vanishes at nominal parameters says nothing about the shift
$\delta q_\theta\Delta\theta$; S06 found shifts of micrometres for
$\Delta\theta$ of order $10^{-3}$.

\begin{remark}[Number of modes]
With $m$ complex modes and one fixed phase there are $2m-1$ real unknowns,
fewer than the pupil observations. Problem \cref{eq:faithful-problem} is then
a least-squares fit whose residual is the faithfulness error itself; an exact
zero is not expected and is not the acceptance test. This is the opposite of
the S04 regime, where more source coordinates than force equations made the
frozen fit exact on the design grid and uninformative off it.
\end{remark}

\subsection{The acoustics-off meniscus as the main shaping element}
Write the displaced volume of interface $j$ as $V_j$ and its pinned rim radius
as $R$. With $\nabla E(q_0)=0$, the acoustic load must supply only the
difference between the target and the meniscus selected by $(V_0,V_1)$:
\begin{equation}
 F_{\mathrm{ac}}(q^\star;c)=\nabla E(q^\star)
 =\nabla E(q^\star)-\nabla E(q_0),
\end{equation}
which vanishes when the target is itself a meniscus. In zero gravity the
meniscus at fixed volume is the spherical cap through the rim, the surface of
constant mean curvature; with gravity it solves the axisymmetric
Young--Laplace equation $\sigma\kappa=\Delta\rho\,g\,h-\lambda_j$ with a
multiplier $\lambda_j$ for the volume \cite{finn1986}. Fluidic fabrication
uses the same principle with density matching and solidification
\cite{elgarisi2021}. Here the meniscus is the starting state of the acoustic
problem, so the fill volumes and the free optical parameters of the target
(for a shared-conjugate pair, the vertex positions at fixed conjugates)
should be chosen first, so that the pupils of the target lie as close as
possible to an acoustics-off meniscus. Gravity enters only through the Bond
number $\Delta\rho\,gR^2/\sigma$. It may be treated as part of the meniscus,
or the problem may be posed first at zero gravity, as with density-matched
fluids, and gravity then corrected as a parasitic load by the acoustics.

This ordering reduces the acoustic task from forming the lens to correcting
the residual asphericity between a meniscus and the optical surface. Because
the first-order shifts $\delta q_\theta\Delta\theta$ scale with the acoustic
load, a smaller correction is also less sensitive to the parameter tolerances
of \cref{eq:faithful-problem}, which is the robustness failure observed when
the whole lens shape is supplied acoustically.

\subsection{Field energy and the scale of force wells}
For a standing wave of wavenumber $k$ the radiation stress on an interface
varies on the length $\pi/k=\lambda/2$, so the equilibria of
\cref{eq:held-equilibrium} at fixed $c$ are separated by displacements of order
$\lambda/4$. A target whose sag $h$ is comparable to $\lambda/4$ lies in one
well among several, and the unforced start may belong to another; at
3.6~MHz in water $\lambda/4\approx104~\mu$m while the S04 vertices move by
$150~\mu$m. Choosing $h\ll\lambda/8$, that is $f\ll c/(8h)$, places the
start-up and the target in one well. This is a guideline from the length
scale of the stress, not a theorem; the branch conditions of
\cref{prop:adiabatic} are the test. Capping the field energy limits the
radiation stress to the scale of the required load, of order
$p_{\max}^2/(\rho c^2)$ times the acoustic contrast, so that the load is not
the small difference of large cancelling stresses that made S04 sensitive
to parameters.

\subsection{Acceptance}
A command is accepted as faithful when, in the declared model,
\begin{enumerate}
 \item the acoustics-off state $q_0$ at the declared fill volumes is a stable
 capillary-gravity equilibrium, and the start-up branch from $q_0$ satisfies the hypotheses of
 \cref{prop:adiabatic}, and a forward start-up trajectory, integrated with a
 stable scheme at refined time steps, reaches the same state;
 \item that state meets the pupil height goal under fixed-command refinement of
 the wave, the element order and the surface space;
 \item the first-order shifts $\delta q_\theta\Delta\theta$ for the declared
 tolerances stay within the goal;
 \item the coupled spectrum at $s=1$ keeps a declared margin $\lambda$.
\end{enumerate}
Optical ray gates are applied to the accepted state. None of these conditions
addresses non-axisymmetric modes, fluid inertia, acoustic streaming or
heating, and none is experimental validation.
